\documentclass[12pt,a4paper]{ctexart}

\usepackage{amsmath,amssymb,bm,booktabs,caption,enumitem,fancyhdr,float,geometry,hyperref,tcolorbox}

\geometry{left=25mm,right=25mm,top=25mm,bottom=25mm}
\hypersetup{colorlinks=true,linkcolor=blue,citecolor=blue,urlcolor=blue}
\setlength{\parindent}{2em}
\setlength{\parskip}{0.35em}
\setlength{\headheight}{15pt}
\pagestyle{fancy}
\fancyhf{}
\lhead{Full Material-Volume EP Flux}
\rhead{\thepage}

\newcommand{\p}{\partial}
\newcommand{\dd}{\,\mathrm d}
\newcommand{\D}{\mathrm D}
\newcommand{\meanV}[1]{\left\langle #1\right\rangle_{\mathcal V}}
\newcommand{\meanL}[1]{\left\langle #1\right\rangle_{\ell}}
\newcommand{\fluc}[1]{#1^{\dagger}}
\newcommand{\flucV}[1]{\left(#1\right)^{\dagger}_{V}}
\newcommand{\flucL}[1]{\left(#1\right)^{\dagger}_{\ell}}
\newcommand{\bx}{\boldsymbol x}
\newcommand{\by}{\boldsymbol y}
\newcommand{\bu}{\boldsymbol u}
\newcommand{\bU}{\boldsymbol U}
\newcommand{\bR}{\boldsymbol R}
\newcommand{\bV}{\boldsymbol V}
\newcommand{\bt}{\boldsymbol t}
\newcommand{\bn}{\boldsymbol n}
\newcommand{\PP}{\boldsymbol{\mathsf P}_{\perp}}
\newcommand{\II}{\boldsymbol{\mathsf I}}
\newcommand{\GG}{\mathcal G}
\newcommand{\LL}{\mathcal L}
\newcommand{\RR}{\mathcal R}
\newcommand{\VV}{\mathcal V_c}
\newcommand{\SSS}{\partial\mathcal V_c}
\newcommand{\MM}{\mathcal M}
\newcommand{\UU}{\mathcal U}
\newcommand{\KK}{\mathcal K}
\newcommand{\SSrc}{\mathcal S}

\title{\bfseries 大曲率 Material-Volume EP Flux 完整推广：\\局地通量、非局地平衡响应与材料体预算}
\author{ZONCA-EP-FLUX-TEST}
\date{2026 年 9 月 9 日}

\begin{document}
\maketitle

本文档是对 \texttt{large\_curvature\_material\_volume\_ep\_flux.tex} 的独立修正版推广。
原有理论转向是正确的：当 $\kappa r\sim 10\text{--}35$ 且 metric valid fraction 的中位数为 $0$ 时，
当前中尺度涡不能再被当作小曲率细涡管处理，必须改用全局 Cartesian 坐标中的 coherent 材料体
$\VV(t)\subset\mathbb R^3$。本修正版处理两个会破坏闭合逻辑的关键点：
\begin{enumerate}[label=\arabic*.]
  \item 非局地 QG/PV 平衡反演不能默认写成局地通量散度；
  \item 全材料体平均量不能再被当作空间场取散度。
\end{enumerate}
本版新增一节，专门说明热/浮力通量与水平动量通量在大曲率材料体中的关系。

\begin{tcolorbox}[colback=blue!3,colframe=blue!65!black,title={修正后的核心口径},fonttitle=\bfseries]
大曲率 Material-Volume EP Flux 的主对象不是一个单一的局地 EP 通量向量，
而是由三部分共同组成的诊断预算：
\[
  \boxed{
  \GG_i(\bx,t)
  =
  -\rho_0^{-1}\p_j\!\left(\rho_0 R_{ij}^{\ell}(\bx,t)\right)
  +
  L_i(\bx,t)
  +
  \RR_i(\bx,t)
  }.
\]
其中 $R_{ij}^{\ell}$ 是局地粗粒化 Reynolds stress 场，可以取散度；
$L_i=\LL_i[B^{\ell},P^{\ell},q_c;\VV,\mathcal D,\mathcal B]$ 是由 QG/Boussinesq PV 反演、
热成风约束和边界条件共同定义的非局地平衡响应；$\RR_i$ 收集年龄转、混合、
非保守 PV、边界识别误差和尺度截断残差。
\end{tcolorbox}

\section{问题诊断与修正原则}

传统 TEM/E-P flux 的紧凑形式
\[
  \p_t\overline u-f_0v^*
  =
  \rho_0^{-1}\nabla\cdot\boldsymbol F+\overline X
\]
成立在明确的平均、坐标和准地转尺度假设下。把它推广到大曲率 coherent 涡旋时，
不能只把纬向平均换成材料体平均、再把二维散度换成三维散度。

第一，全体平均
\[
  \meanV{\phi}(t)
  =
  |\VV(t)|^{-1}\int_{\VV(t)}\phi(\bx,t)\,\dd V
\]
是材料体上的整体统计量。它没有内部空间坐标依赖，因此 $\p_j\meanV{\phi}$ 不是 coherent 体内部的
局地散度。它只能用于总预算、质心速度、生命周期合成和统计检验。

第二，QG/PV 反演是边值问题。给定 PV、浮力通量或热成风约束后的速度响应依赖整个区域的源项、
边界条件、背景稳定度和 $f_0,\beta$ 等参数。除非已经构造并验证局地化近似核，否则不能把该响应写成
某个未验证的局地映射通量散度
\[
  \p_j\left(\rho_0 H_{ij}^{\mathrm{loc}}\right)
\]
并进一步用 Gauss 定理变成边界通量。正确做法是把它保留为非局地响应算子。

\section{材料体、全体平均与局地粗粒化场}

令 coherent 材料体为
\[
  \VV(t)=
  \{\bx:\ q_c(\bx,t)\in Q_c,\ \bx\ \hbox{属于同一 coherent 连通核}\}.
\]
边界 $\SSS(t)$ 可以是严格材料面，也可以是由离散阈值、生命周期合成或目标追踪算法给出的诊断边界。
若边界不是严格材料面，穿透误差必须保留为显式残差。

全体平均定义为
\[
  \meanV{\phi}(t)
  =
  \frac{1}{|\VV(t)|}\int_{\VV(t)}\phi(\bx,t)\,\dd V,
  \qquad
  \flucV{\phi}(\bx,t)=\phi(\bx,t)-\meanV{\phi}(t).
\]
它适用于整体量：
\[
  M_q(t)=\int_{\VV(t)}q_c\,\dd V,\qquad
  \bR_c(t)=M_q^{-1}\int_{\VV(t)}\bx q_c\,\dd V.
\]
这些量不构成材料体内部的空间通量场，因此不进入 $\p_j(\cdot)$。

为了定义可取散度的 EP 型局地强迫，引入局地粗粒化平均：
\[
  \meanL{\phi}(\bx,t)
  =
  \frac{\int_{\VV(t)}W_\ell(\bx,\by)\phi(\by,t)\,\dd V_{\by}}
       {\int_{\VV(t)}W_\ell(\bx,\by)\,\dd V_{\by}},
  \qquad
  \flucL{\phi}(\bx,t)=\phi(\bx,t)-\meanL{\phi}(\bx,t).
\]
只有依赖 $\bx$ 的 $\meanL{\cdot}$ 才能生成可取散度的局地通量场。

\section{Cartesian 局地通量场}

在全局 Cartesian 坐标 $x_i=(x,y,z)$ 中定义局地粗粒化通量：
\[
  R_{ij}^{\ell}(\bx,t)
  =
  \meanL{\flucL{u_i}\flucL{u_j}},
  \qquad
  B_i^{\ell}(\bx,t)
  =
  \meanL{\flucL{u_i}\flucL{b}},
  \qquad
  P_i^{\ell}(\bx,t)
  =
  \meanL{\flucL{u_i}\flucL{q}}.
\]
这里 $R_{ij}^{\ell}$、$B_i^{\ell}$、$P_i^{\ell}$ 是材料体内的局地场，而不是整个材料体上的单一数值。
因此应力散度
\[
  A_i^{R}(\bx,t)
  =
  -\rho_0^{-1}\p_j\!\left(\rho_0 R_{ij}^{\ell}(\bx,t)\right)
\]
是合法的局地动量强迫诊断。

\section{热通量与水平动量通量的关系}

先把被比较的对象写清楚。水平动量通量是 Reynolds stress 的水平分量，例如
\[
  R_{xy}^{\ell}=\meanL{\flucL{u}\flucL{v}},
  \qquad
  R_{xz}^{\ell}=\meanL{\flucL{u}\flucL{w}},
\]
其中 $R_{xy}^{\ell}$ 表示扰动经向输送纬向动量，$R_{xz}^{\ell}$ 表示扰动垂向输送纬向动量。
热通量在 Boussinesq/QG 口径下写成浮力通量
\[
  B_y^{\ell}=\meanL{\flucL{v}\flucL{b}},
  \qquad
  B_z^{\ell}=\meanL{\flucL{w}\flucL{b}}.
\]
若使用位温 $\theta$，可通过 $b=g\theta'/\theta_0$ 或等价线性状态方程把 $\theta$ 通量换成
浮力通量。

在传统近直轴 TEM 极限中，二者可以组成局地 EP 通量：
\[
  F_y=-\rho_0\overline{u'v'},
  \qquad
  F_z=\rho_0\frac{f_0}{N^2}\overline{v'b'}.
\]
于是平均纬向动量强迫为
\[
  \rho_0^{-1}\nabla\cdot\boldsymbol F
  =
  -\rho_0^{-1}\p_y\left(\rho_0\overline{u'v'}\right)
  +
  \rho_0^{-1}\p_z
  \left(\rho_0\frac{f_0}{N^2}\overline{v'b'}\right).
\]
这说明传统 EP flux 中热通量不是普通垂直动量通量，而是经由热成风和平衡关系折算成
对纬向动量的垂向波活动通量分量。

在大曲率材料体中，这个局地相加关系不再是基本闭合式。水平动量通量仍然给出局地应力散度：
\[
  A_x^{R}
  =
  -\rho_0^{-1}
  \left[
    \p_x(\rho_0R_{xx}^{\ell})
    +
    \p_y(\rho_0R_{xy}^{\ell})
    +
    \p_z(\rho_0R_{xz}^{\ell})
  \right].
\]
其中与水平动量输送最直接相关的是 $\p_y(\rho_0R_{xy}^{\ell})$ 和 $\p_z(\rho_0R_{xz}^{\ell})$。

热/浮力通量的作用则进入非局地平衡响应。形式上令
\[
  \MM[\psi^{\mathrm{resp}}]
  =
  \SSrc_B[B^{\ell}]
  +
  \SSrc_P[P^{\ell}]
  +
  \SSrc_c[q_c]
  \quad \hbox{in } \mathcal D,
  \qquad
  \psi^{\mathrm{resp}}\big|_{\partial\mathcal D}\in\mathcal B.
\]
解为
\[
  \psi^{\mathrm{resp}}
  =
  \MM^{-1}
  \left(
    \SSrc_B[B^{\ell}]
    +
    \SSrc_P[P^{\ell}]
    +
    \SSrc_c[q_c]
  \right).
\]
由流函数到速度的平衡映射记为 $\UU_i$，则热通量对第 $i$ 个动量或速度响应分量的贡献为
\[
  L_i^{B}(\bx,t)
  =
  \UU_i\MM^{-1}\SSrc_B[B^{\ell}](\bx,t).
\]
PV 通量和 coherent PV anomaly 的响应类似：
\[
  L_i^{P}=\UU_i\MM^{-1}\SSrc_P[P^{\ell}],
  \qquad
  L_i^{c}=\UU_i\MM^{-1}\SSrc_c[q_c],
\]
因此
\[
  L_i=L_i^B+L_i^P+L_i^c .
\]

于是水平动量通量与热通量在大曲率材料体中的关系是
\[
  \boxed{
  \GG_x
  =
  A_x^R
  +
  L_x^B
  +
  L_x^P
  +
  L_x^c
  +
  \RR_x
  }.
\]
它们仍然共同决定同一个平均/材料体强迫，但关系不再是
``一个水平动量通量项加一个热通量局地散度项''，而是
\[
  \boxed{
  \begin{aligned}
  \hbox{水平动量通量}
  &\longrightarrow
  A_x^R\ \hbox{局地散度},\\
  \hbox{热/浮力通量}
  &\longrightarrow
  L_x^B=\UU_x\MM^{-1}\SSrc_B[B^{\ell}]\ \hbox{非局地响应}.
  \end{aligned}
  }
\]

若要恢复类似传统 EP flux 的局地关系，必须额外满足尺度分离和局地化条件。即存在局地响应通量
$H_{ij}^B$，使得
\[
  L_i^B
  =
  \rho_0^{-1}\p_j(\rho_0H_{ij}^{B})
  +
  \epsilon_i^{B},
  \qquad
  |\epsilon_i^{B}|\ll |L_i^B|.
\]
只有此时才可把热通量写回某种 EP 型局地通量分量。否则最诚实的关系就是非局地算子关系：
\[
  B^{\ell}
  \ \xrightarrow{\ \SSrc_B\ }\ 
  \hbox{PV/thermal source}
  \ \xrightarrow{\ \MM^{-1}\ }\ 
  \psi^{\mathrm{resp}}
  \ \xrightarrow{\ \UU_i\ }\ 
  L_i^B .
\]

物理上，水平动量通量直接描述扰动把水平动量搬到哪里；热/浮力通量先改变密度和热成风剪切，
再通过 PV 反演调整速度场和 PV 分布。两者不是孤立诊断量，而是在
\[
  \GG_i=A_i^R+L_i+\RR_i
\]
中共同闭合到材料体强迫。

\section{非局地 QG/PV 平衡响应算子}

令 $\mathcal D$ 为包含 $\VV(t)$ 的计算域，$\mathcal B$ 为边界条件集合。QG/Boussinesq 层级下，
平衡响应可抽象为如下边值问题：
\[
  \MM[\psi^{\mathrm{resp}}]
  =
  \SSrc\!\left[B^{\ell},P^{\ell},q_c,\overline q,N^2,f_0,\beta\right]
  \quad \hbox{in } \mathcal D,
  \qquad
  \psi^{\mathrm{resp}}\big|_{\partial\mathcal D}\in\mathcal B.
\]
速度响应写为
\[
  u_i^{\mathrm{resp}}(\bx,t)
  =
  \UU_i[\psi^{\mathrm{resp}}](\bx,t),
\]
于是非局地平衡强迫定义为
\[
  L_i(\bx,t)
  =
  \LL_i\!\left[B^{\ell},P^{\ell},q_c;\VV,\mathcal D,\mathcal B\right](\bx,t).
\]
$L_i$ 的值依赖整个源场和边界条件，而不只依赖 $\bx$ 附近的 $B_i^{\ell}$。因此完整
Material-Volume EP 型局地诊断为
\[
  \boxed{
  \GG_i(\bx,t)
  =
  A_i^{R}(\bx,t)
  +
  L_i(\bx,t)
  +
  \RR_i(\bx,t)
  }.
\]
其中 $\RR_i$ 至少包含年龄转或非 QG 残差、混合和摩擦、非保守 PV 源、粗粒化截断误差、
非材料边界穿透误差，以及三维 coherent 体投影成分层或壳层时的几何误差。

\section{正确的体积分与边界积分预算}

对材料体积分得到总强迫预算：
\[
  \int_{\VV}\GG_i\,\dd V
  =
  \int_{\VV}A_i^{R}\,\dd V
  +
  \int_{\VV}L_i\,\dd V
  +
  \int_{\VV}\RR_i\,\dd V.
\]
其中只有 Reynolds stress 散度项可直接应用 Gauss 定理：
\[
  \int_{\VV}A_i^{R}\,\dd V
  =
  -\int_{\SSS}R_{ij}^{\ell}n_j\,\dd S
  -
  \int_{\VV}R_{ij}^{\ell}\p_j(\ln\rho_0)\,\dd V.
\]
若 $\rho_0$ 在所用坐标中近似常数，则
\[
  \int_{\VV}A_i^{R}\,\dd V
  \simeq
  -\int_{\SSS}R_{ij}^{\ell}n_j\,\dd S.
\]
平衡响应项必须保留为 $\int_{\VV}L_i\,\dd V$，除非独立构造并验证局地通量 $H_{ij}$。此时才可写
\[
  \int_{\VV}L_i\,\dd V
  =
  \int_{\SSS}H_{ij}n_j\,\dd S
  +
  \int_{\VV}\epsilon_i^{H}\,\dd V.
\]
若边界不是严格材料面，还需显式保留穿透项
\[
  \mathcal E_{\partial V}[\phi]
  =
  \int_{\SSS}\phi\left(\bu-\bU_{\SSS}\right)\cdot\bn\,\dd S.
\]

\section{PV 质心速度与 coherent tilt/bend 诊断}

PV 质心定义为
\[
  \bR_c(t)=
  \frac{\int_{\VV}\bx q_c\,\dd V}{\int_{\VV}q_c\,\dd V}.
\]
对其求导得到
\[
  \bV_c
  =
  \frac{1}{M_q}\int_{\VV}\bu q_c\,\dd V
  +
  \frac{1}{M_q}\int_{\VV}(\bx-\bR_c)S_q\,\dd V
  +
  \bV_{\partial V}.
\]
分解为
\[
  V_{c,i}
  =
  \meanV{u_i}
  +
  \frac{|\VV|}{M_q}\meanV{\flucV{u_i}\flucV{q_c}}
  +
  V_{\mathrm{src},i}
  +
  V_{\partial V,i}.
\]
这条式子使用的是全体平均，因而只用于质心速度本身，不对其内部通量项取空间散度。

若用分层 PV 质心构造 axis skeleton，令弧长参数为 $s$，切向量为
\[
  \bt=\p_s\bR_c(s,t),\qquad |\bt|=1.
\]
轴线形态演化满足
\[
  \D_t\bt
  =
  \p_s\bV_c
  -
  \left(\bt\cdot\p_s\bV_c\right)\bt
  =
  \PP\,\p_s\bV_c,
  \qquad
  \PP=\II-\bt\bt.
\]
因此 $ \PP\,\p_s\bV_c $ 是形态响应，不是 EP 强迫本身；EP/PV 诊断通过 $\bV_c$ 的非局地响应进入该式。

\section{传统 EP flux 与 curved-tube 的极限对应}

当以下条件同时满足时，大曲率材料体框架可退化到 curved-tube 或传统 TEM/E-P flux：
\begin{enumerate}[label=\arabic*.]
  \item coherent 体可由单一小曲率细管坐标覆盖，$\kappa r\ll1$ 且 $J=1-\kappa_\alpha x^\alpha>0$；
  \item 局地截面平均可稳定替代全局 Cartesian 粗粒化；
  \item 平衡响应算子 $L_i$ 在所研究尺度上可局地化，且局地化残差已被量化；
  \item 只关心一个平均动量分量，并且几何联络项在小曲率展开内可控。
\end{enumerate}
此时可近似恢复
\[
  \nabla_aF_{\mathrm{CT}}^{ai}
  =
  J^{-1}\p_a(JF_{\mathrm{CT}}^{ai})
  +
  \Gamma^i{}_{ka}F_{\mathrm{CT}}^{ak}.
\]
但对于当前验证样本，$\kappa r$ 的中位数约为 $10$ 到 $35$，metric valid fraction 的中位数为 $0$。
因此 curved-tube 几何项只能作为尺度风险审计，不能作为已经闭合的动力 forcing。

\section{数值实现路线和验证标准}

第一版数值实现应按以下顺序推进：
\begin{enumerate}[label=\arabic*.]
  \item 用 coherent PV anomaly 或流函数核识别三维连通材料体 $\VV(t)$；
  \item 同时输出全体平均量和局地粗粒化场，不混用二者；
  \item 在 Cartesian 网格或分层壳上计算 $R_{ij}^{\ell}$、$B_i^{\ell}$、$P_i^{\ell}$；
  \item 只对 $R_{ij}^{\ell}$ 的局地散度计算 $A_i^R$；
  \item 使用 QG/Boussinesq PV 反演求解非局地响应 $L_i$，并记录边界条件与源项定义；
  \item 对 $\GG_i=A_i^R+L_i+\RR_i$ 做体积分预算，边界通量只用于已经散度化的部分；
  \item 计算 PV 质心速度 $\bV_c$ 与 $\PP\,\p_s\bV_c$，诊断 tilt、bend 和 coherent 失稳；
  \item 用 track-level accumulator 做 bootstrap/jackknife，避免从最终均值代表涡伪造置信区间。
\end{enumerate}

最低验证标准为：
\begin{enumerate}[label=\arabic*.]
  \item 局地场 $R_{ij}^{\ell},B_i^{\ell},P_i^{\ell}$ 与全体平均 $\meanV{\cdot}$ 在代码和报告中分开命名；
  \item 不出现把浮力响应默认写成局地散度闭合的公式；
  \item 不对全体平均 Reynolds stress 取空间散度；
  \item 报告 $\int_{\VV}L_i\,\dd V$、$\int_{\VV}\RR_i\,\dd V$ 和边界穿透项的量级；
  \item curved-tube 项只在小曲率审计通过时进入动力解释。
\end{enumerate}

\section{一句话总结}

大曲率 coherent 涡旋的 EP-flux 推广不是把传统 EP 向量机械改写成三维局地散度，而是把可散度化的
局地 Reynolds stress、非局地 QG/PV 平衡响应、PV 质心动力和真实边界残差放进同一个材料体预算。
热/浮力通量与水平动量通量仍然共同决定同一个材料体强迫，但前者主要通过非局地平衡反演进入，
后者则通过局地 Reynolds stress 散度进入。这样既保留了 E-P 理论中合成解释的核心精神，也避免了
把非局地反演误写成局地边界通量、以及对全体平均量取空间散度这两个闭合错误。

\end{document}
